\documentclass[11pt]{article}
\usepackage[T1]{fontenc}
\usepackage{lmodern,amsmath,amssymb,amsthm,mathtools}
\usepackage[margin=1in]{geometry}
\usepackage{microtype,booktabs,longtable,array,enumitem,xcolor,hyperref}
\hypersetup{colorlinks=true,linkcolor=blue!45!black,citecolor=blue!45!black,
 urlcolor=blue!45!black,pdftitle={Linear forms, moment determinants, and arithmetic constraints for even beta values},pdfauthor={twinheterodox}}
\urlstyle{same}
\setlength{\emergencystretch}{2em}
\setlist{nosep}
\newtheorem{theorem}{Theorem}[section]
\newtheorem{proposition}[theorem]{Proposition}
\newtheorem{lemma}[theorem]{Lemma}
\newtheorem{corollary}[theorem]{Corollary}
\theoremstyle{remark}\newtheorem{remark}[theorem]{Remark}
\numberwithin{equation}{section}
\newcommand{\Q}{\mathbb Q}\newcommand{\Z}{\mathbb Z}
\newcommand{\R}{\mathbb R}\newcommand{\C}{\mathbb C}
\newcommand{\dd}{\,\mathrm d}\newcommand{\E}{\mathbb E}
\newcommand{\bet}{\beta(4)}\newcommand{\zet}{\zeta(4)}
\DeclareMathOperator{\rank}{rank}\DeclareMathOperator{\diag}{diag}
\DeclareMathOperator{\lcm}{lcm}\DeclareMathOperator{\cont}{cont}
\DeclareMathOperator{\Span}{span}
\title{Linear forms, moment determinants, and arithmetic constraints for even beta values}
\author{twinheterodox}
\date{13 November 2026}
\begin{document}
\maketitle
\begin{abstract}
This supplement studies arithmetic constraints in attempts to prove
irrationality of an individual even Dirichlet beta value or of a member
of a small explicit set. For a positive moment construction involving
$\beta(4)$ and $\zeta(4)$, we determine the exact power of two in every
principal minor of its beta-coefficient matrix. The resulting necessary
denominator growth rules out a specified certificate based on one-sided
rational approximation to $\zeta(4)$. We also give controlled rational
elimination of the unwanted period, uniform divisibility results for
compressed determinants, and a two-sided estimate for the residual decay
of a thirteen-value Gram construction. A separate linear-form construction
provides eleven independent, uniformly small integer forms in fourteen
candidate beta values, with nonvanishing still unresolved. These results
identify usable mechanisms and precise remaining obstacles; they do not
establish a new irrationality theorem.
\end{abstract}
\setcounter{tocdepth}{1}
\begingroup\small
\tableofcontents
\endgroup
\clearpage
\input{paper/sections/01-introduction.tex}
\input{paper/sections/02-moments.tex}
\input{paper/sections/03-valuations.tex}
\input{paper/sections/04-substitution.tex}
\input{paper/sections/05-compression.tex}
\input{paper/sections/06-content.tex}
\input{paper/sections/07-gram.tex}
\input{paper/sections/08-linear.tex}
\input{paper/sections/09-conclusions.tex}
\appendix
\input{paper/sections/10-appendix.tex}
\input{paper/references.tex}
\end{document}
